\documentclass[../DoAn.tex]{subfiles}
\begin{document}

\section{Đặt vấn đề}
\label{sec:dvd}

Trong bối cảnh chuyển đổi số đang diễn ra mạnh mẽ trong lĩnh vực giáo dục, việc ứng dụng công nghệ thông tin vào công tác quản lý và vận hành các trung tâm đào tạo đã trở thành yêu cầu tất yếu. Đối với các trung tâm giảng dạy ngoại ngữ, ngoài hoạt động chuyên môn là giảng dạy và đào tạo, công tác quản lý tài sản giữ vai trò quan trọng trong việc đảm bảo hoạt động ổn định, tiết kiệm chi phí và nâng cao chất lượng dịch vụ.

Tài sản của một trung tâm giảng dạy ngoại ngữ bao gồm nhiều loại khác nhau như phòng học, thiết bị giảng dạy, máy tính, máy chiếu, hệ thống âm thanh và các cơ sở vật chất hỗ trợ đào tạo. Cùng với quá trình mở rộng quy mô và nâng cao chất lượng đào tạo, số lượng và giá trị tài sản của trung tâm ngày càng tăng, kéo theo yêu cầu quản lý ngày càng chặt chẽ và chính xác.

Trên thực tế, tại nhiều trung tâm giảng dạy ngoại ngữ, công tác quản lý tài sản vẫn chủ yếu được thực hiện bằng các phương pháp thủ công hoặc thông qua các công cụ đơn giản như bảng tính. Cách tiếp cận này tồn tại nhiều hạn chế, bao gồm việc khó theo dõi tình trạng và vòng đời tài sản, dễ xảy ra sai sót trong quá trình cập nhật dữ liệu, thiếu tính đồng bộ và không hỗ trợ hiệu quả cho công tác thống kê, báo cáo cũng như ra quyết định quản lý.

Xuất phát từ những bất cập nêu trên, việc xây dựng một hệ thống quản lý tài sản dựa trên công nghệ thông tin, phù hợp với đặc thù của trung tâm giảng dạy ngoại ngữ, là hết sức cần thiết. Đề tài đồ án này được thực hiện nhằm nghiên cứu và triển khai một giải pháp quản lý tài sản góp phần nâng cao hiệu quả quản lý, tối ưu hóa việc sử dụng tài nguyên và hỗ trợ công tác vận hành của trung tâm.

\section{Các giải pháp hiện tại và hạn chế}
\label{sec:giaiphap}

Hiện nay, các giải pháp quản lý tài sản đang được áp dụng trong các tổ chức giáo dục có thể chia thành hai nhóm chính. Nhóm thứ nhất là các phương pháp quản lý thủ công hoặc bán tự động, bao gồm việc sử dụng sổ sách, biểu mẫu hoặc các bảng tính như Excel. Các phương pháp này có ưu điểm là dễ triển khai, chi phí thấp và phù hợp với các đơn vị có quy mô nhỏ. Tuy nhiên, khi số lượng tài sản tăng lên, những hạn chế của nhóm giải pháp này ngày càng rõ rệt, như khó đảm bảo tính nhất quán của dữ liệu, không hỗ trợ theo dõi lịch sử sử dụng và bảo trì tài sản, cũng như tiềm ẩn nguy cơ thất lạc hoặc sai lệch thông tin.

Nhóm thứ hai là các hệ thống phần mềm quản lý tài sản chuyên dụng. Các hệ thống này thường cung cấp nhiều chức năng như quản lý vòng đời tài sản, phân công tài sản cho người sử dụng hoặc phòng ban, theo dõi tình trạng tài sản và hỗ trợ báo cáo thống kê. Mặc dù mang lại hiệu quả cao, nhiều giải pháp thương mại có chi phí triển khai và vận hành lớn, yêu cầu hạ tầng kỹ thuật phức tạp hoặc chưa được tùy biến phù hợp với mô hình hoạt động của các trung tâm giảng dạy ngoại ngữ tại Việt Nam.

Từ những phân tích trên có thể thấy rằng, các giải pháp hiện tại vẫn chưa đáp ứng đầy đủ yêu cầu về tính linh hoạt, chi phí hợp lý và khả năng triển khai thực tế đối với các trung tâm giảng dạy ngoại ngữ có quy mô vừa và nhỏ. Do đó, việc nghiên cứu và lựa chọn một giải pháp phù hợp hơn là cần thiết.

\section{Mục tiêu và định hướng giải pháp}

Mục tiêu của đồ án là xây dựng một hệ thống quản lý tài sản dành cho trung tâm giảng dạy ngoại ngữ, nhằm hỗ trợ quản lý tập trung và theo dõi chính xác các thông tin liên quan đến tài sản. Hệ thống hướng tới việc chuẩn hóa quy trình quản lý, giảm thiểu sai sót trong quá trình vận hành và nâng cao hiệu quả thống kê, báo cáo phục vụ công tác quản lý.

Để đạt được mục tiêu trên, đồ án định hướng triển khai giải pháp dựa trên nền tảng phần mềm quản lý tài sản mã nguồn mở ứng dụng quản lý tài sản. Trên cơ sở nền tảng này, hệ thống sẽ được cấu hình và điều chỉnh để phù hợp với đặc thù của trung tâm giảng dạy ngoại ngữ, bao gồm quản lý danh mục tài sản, phân loại theo phòng học và mục đích sử dụng, theo dõi tình trạng và lịch sử sử dụng tài sản, cũng như phân quyền người dùng. Định hướng này cho phép tận dụng các ưu điểm sẵn có của nền tảng, đồng thời giảm chi phí và thời gian triển khai hệ thống.

\section{Đóng góp của đồ án}

Đồ án tập trung vào việc nghiên cứu và triển khai một hệ thống quản lý tài sản có tính ứng dụng thực tiễn cao. Các đóng góp chính của đồ án bao gồm việc phân tích bài toán quản lý tài sản tại trung tâm giảng dạy ngoại ngữ, đề xuất mô hình quản lý phù hợp và triển khai thử nghiệm hệ thống quản lý tài sản. Ngoài ra, đồ án còn xây dựng tài liệu hướng dẫn triển khai và sử dụng hệ thống, làm cơ sở tham khảo cho việc áp dụng và mở rộng trong các trung tâm đào tạo tương tự.

\section{Bố cục đồ án}

Phần còn lại của báo cáo đồ án tốt nghiệp được tổ chức như sau. Chương 2 trình bày cơ sở lý thuyết và các khái niệm liên quan đến quản lý tài sản, bao gồm quản lý vòng đời tài sản và các mô hình quản lý trong tổ chức giáo dục. Chương 3 phân tích yêu cầu hệ thống và đề xuất kiến trúc tổng thể cho hệ thống quản lý tài sản của trung tâm giảng dạy ngoại ngữ. Chương 4 mô tả chi tiết quá trình triển khai hệ thống và các chức năng chính. Cuối cùng, Chương 5 trình bày kết quả đạt được, đánh giá hệ thống và đề xuất các hướng phát triển trong tương lai.

\end{document}
